\chapter{Dark matter constraints from the Galactic Center}

The Galactic Center is expected to be the single brightest source of dark matter annihilation products on the sky. In particular, using a standard NFW profile, shown in \cref{NFW}, with $r_s = 20$
kpc and a local density of $0.4~\text{GeV/cm}^3$ \cite{hooper2018tasi}, from \cref{eq:AnnihilationFlux_general} we obtain a flux of gamma rays that is more than three orders of magnitude larger than that predicted
from the brightest dwarf galaxies.
However, the mass within the innermost kpc of the Milky Way is dominated by baryons, 
this leads to a higher astrophysical background and, most importantly, numerical simulations need to include hydrodynamics
and feedback physics in addition to the gravitational effects, resulting in uncertainties in the obtained DM profiles towards the GC.
These simulations agree that the DM profile of MW-like galaxies develops a peak towards the Galactic Center but their resolution prevents them from determining a well-defined dark matter distribution.

Therefore, in this chapter we define a set of representative DM density profiles, determine their parameters from observational constraints and compute their $J$-factors.
Then, we introduce a simple model of thermal relic DM and confront it with the upper limits on the annihilation cross section obtained by the H.E.S.S. Inner Galaxy Survey,
showing how different profiles lead to different constraints.


\section{Dark matter density profiles and $J$-factors} \label{GC}

As long as a clearer determination of the DM profile in the inner Milky Way regions is not achieved, it is worth considering a range of possible choices and compare them. 
In particular, the two limiting cases are: a peaked profile, in which the density keeps rising steeply towards the GC, and a cored one, in which the density flattens to a finite central value. 
Therefore, in what follows we consider the Einasto (peaked) \cite{Einasto1965}, Burkert (cored) \cite{Burkert1995} 
and NFW \cite{NFW1996,NFW1997} with a variable core radius $r_c$ as representative DM density profiles 
$\rho_\mathrm{DM}(r)$.


For the Einasto profile the functional form reads

\begin{equation}
\label{eq:rho_einasto}
\rho_{\mathrm{Ein}}(r)
=
\rho_s \exp\!\left[
-\frac{2}{\alpha}\left(\left(\frac{r}{r_s}\right)^{\alpha}-1\right)
\right],
\end{equation}

where $\alpha = 0.17$ is a shape parameter that controls how the logarithmic slope of the density varies with the radius. Note that the Einasto profile does not diverge for small radii, 
nevertheless, in the region of our analysis it behaves as a peaked profile.

The Burkert profile is

\begin{equation}
\label{eq:rho_burkert}
\rho_{\mathrm{Bur}}(r)
=
\frac{\rho_s}{\left(1+\dfrac{r}{r_s}\right)\left(1+\left(\dfrac{r}{r_s}\right)^2\right)},
\end{equation}

which flattens to $\rho_s$ at small radii, defining a core of size $\sim r_s$.
 

The functional form of the cored NFW is


\begin{equation}
\label{eq:rho_nfw_cored}
\rho_{\mathrm{NFW}}(r,r_c)=
\begin{cases}
\rho_s\,\dfrac{r_s}{r_c}\,\dfrac{1}{\left(1+\dfrac{r_c}{r_s}\right)^2}
& \text{for } r \le r_c \\[10pt]
\rho_s\,\dfrac{r_s}{r}\,\dfrac{1}{\left(1+\dfrac{r}{r_s}\right)^2}
& \text{for } r > r_c
\end{cases}
\end{equation}

Once $r_c$ is fixed, all three profiles depend on two free parameters, $\rho_s$ and $r_s$.
We determine them by fitting the DM density at the Sun's position: $\rho(r_\odot = 8.33~\text{kpc}) = 0.4~\text{GeV/cm}^3$, and the total DM mass within 60~kpc: $M(<60~\text{kpc}) = 4.7 \times 10^{11}\,M_\odot$.
For the NFW profile these conditions yield $r_s \simeq 14.0$~kpc and
$\rho_s \simeq 0.60~\mathrm{GeV/cm^3}$, these values will be adopted throughout this work.

In order to make a more quantitative description, we compute the corresponding $J$-factors, as done in \cite{Lefranc_2016}. 
From now on, we adopt the dimensional definition of the $J$-factor in order to properly compare our results with the experimental values.
The main difference with \cite{Lefranc_2016} is the DM density at the Sun's position: they adopt \(\rho(r_\odot)=0.3~\mathrm{GeV/cm^3}\), while here we use \(\rho(r_\odot)=0.4~\mathrm{GeV/cm^3}\), since recent studies consider this value more accurate. Improved determinations of the local DM density are being carried out, see \cite{Read_2014} for instance.
As a check, we verified that setting \(\rho(r_\odot)=0.3~\mathrm{GeV/cm^3}\) reproduces the values reported in \cite{Lefranc_2016}.
We integrate the $J$-factors over an angular region, centered on the GC, defined as 

\begin{equation}
\label{eq:Omega_obs}
    \Omega_\mathrm{obs} = \Omega(\theta < 1^\circ,\; |b| > 0.3^\circ),
\end{equation}

where $\theta > 0^\circ$ is the polar angle with respect to the GC direction 
and $b$ is the galactic latitude with respect to the disk plane. At the Sun's 
distance $r_\odot = 8.33$~kpc, the angles $1^\circ$ and $0.3^\circ$ correspond 
to projected distances of approximately 145~pc and 44~pc from the GC, respectively. 
The choice of $|b| > 0.3^\circ$ is motivated by the need to mask the large photon backgrounds from the galactic ridge.
The results are reported in \cref{tab:gc_jfactors}, where we use three illustrative values of $r_c$ for the cored NFW.



\begin{table}[H]
\centering
\begin{tabular}{c|c|c|c|c|c}
\hline
 & Einasto & Burkert & $\mathrm{NFW}_{30\,\mathrm{pc}}$ & $\mathrm{NFW}_{300\,\mathrm{pc}}$ & $\mathrm{NFW}_{3\,\mathrm{kpc}}$ \\
\hline
$\rho_s~[\mathrm{GeV/cm^3}]$
& 0.171 & 1.283 & 0.604 & 0.604 & 0.586 \\
\hline
$r_s~[\mathrm{kpc}]$
& 12.9 & 9.7 & 14.0 & 14.0 & 14.3 \\
\hline
$\log_{10}\!\left(J_{\Omega_{\mathrm{obs}}}\,\left[\dfrac{\mathrm{GeV}^2}{\mathrm{cm}^5}\right]\right)$
& 21.96 & 19.35 & 21.57 & 21.16 & 19.77 \\
\hline
\end{tabular}
\caption{\emph{Values of the parameters $\rho_s$ and $r_s$ obtained from the
fitting conditions, and $J$-factors for DM annihilation integrated over the
region $\Omega_{\mathrm{obs}}$ of \cref{eq:Omega_obs}, centered on the
Galactic Center, for the Einasto, Burkert and cored NFW density profiles,
with the NFW core radii $r_c$ reported as subscripts.}}
\label{tab:gc_jfactors}
\end{table}

The values in \cref{tab:gc_jfactors} span more than two orders of magnitude 
between the Einasto profile, the most peaked one, and Burkert, the most cored one.
This has a direct consequence for 
the upper bounds on the DM annihilation cross section: since the predicted flux scales linearly with the
$J$-factor, the upper bound on $\langle \sigma v \rangle$ is inversely 
proportional to it. Consequently, a more peaked profile yields a stronger constraint, namely a lower upper limit on the cross section, than a cored one.
In particular, the upper bound on $\langle \sigma v \rangle$ derived under the Burkert assumption 
is roughly a factor $\sim 4 \times 10^2$ weaker than the one obtained under Einasto.

This shows that, when the analysis is restricted to regions very close to the Galactic Center, the uncertainty on the DM density profile has a very large impact on the constraints on the annihilation cross section.
Instead, if a larger region of observation is considered, the impact of this uncertainty is reduced, since the different profiles show a similar behavior in outer regions.
On the other hand, restricting the region of observation can reduce the astrophysical backgrounds.
This is the reason why small regions close to the GC are particularly interesting for indirect searches, in particular for the study of DM spikes.


\section{Thermal relic cross section and indirect detection constraints}


We now consider the model where two DM particles annihilate into two BSM mediators, i.e. particles that are not SM, 

\begin{equation}
    \chi\chi \to \phi\phi .
\end{equation}

The mediators are then assumed to later decay into SM particles.

In the following calculation we use natural units, where \(\hbar=c=1\), masses and energies are expressed in GeV and cross sections in $\mathrm{GeV}^{-2}$. 


The annihilation cross section for the process \(\chi\chi\to\phi\phi\) is

\begin{equation}
\label{eq:sigmav_dark_sector}
\langle \sigma v \rangle_{\chi\chi\to\phi\phi}
=
\frac{\pi\alpha_D^2}{m_{\mathrm{DM}}^2},
\end{equation}

where \(m_{\mathrm{DM}}\) is the DM mass and \(\alpha_D\) is the coupling constant of 
the dark sector.

Assuming the thermal relic scenario, the observed DM abundance is obtained for an annihilation cross section of order \(10^{-26}\,\mathrm{cm^3/s}\).
For a quantitative analysis, we use the value for self-conjugate dark matter with mass greater than 10 GeV \cite{Steigman_2012}

\begin{equation}
  \langle\sigma v\rangle =  2.2\times10^{-26}~\mathrm{cm^3/s}
    \simeq
    1.89\times10^{-9}~\mathrm{GeV}^{-2}.
\end{equation}

From \cref{eq:sigmav_dark_sector} we can easily compute the value of $\alpha_D(m_{\mathrm{DM}})$ that gives the correct relic abundance

\begin{equation}
\label{eq:alphaD_solution}
\alpha_D(m_{\mathrm{DM}})
\simeq
2.45\times10^{-5}
\left(
\frac{m_{\mathrm{DM}}}{\mathrm{GeV}}
\right)
\simeq
0.0245
\left(
\frac{m_{\mathrm{DM}}}{1~\mathrm{TeV}}
\right).
\end{equation}

Once \(\alpha_D\) is chosen, the annihilation cross section of our model is fixed.

We now assume that the mediators decay into $W$ bosons, therefore the complete process is 

\begin{equation}
    \chi\chi\to\phi\phi\to W^{+}W^{-}\,W^{+}W^{-} .
\end{equation}


For simplicity, we use the limit for DM annihilating directly into the $W^{+}W^{-}$ channel. 

This is a good approximation because, as discussed in \cref{Gammarays}, the $\gamma$-ray spectrum from $W$ 
hadronization is a broad continuum: the extra step through the mediators halves the average energy of each $W$ and 
doubles their multiplicity, but this mildly affects the spectrum. The resulting limits agree with the 
direct-channel ones up to $\mathcal{O}(1)$ factors.

Ref.~\cite{Abdalla_2022} provides 95\% confidence level (C.L.) upper limits on the annihilation cross section $\langle\sigma v\rangle_{\mathrm{lim}}(m_{\mathrm{DM}})$ as a function of the DM mass. 
H.E.S.S. adopts an NFW profile with $\rho_\odot = 0.39~\mathrm{GeV/cm^3}$ at $r_\odot = 8.5$~kpc and $r_s = 21$~kpc, while we adopt the parameters determined in \cref{GC}: $\rho_\odot = 0.4~\mathrm{GeV/cm^3}$ at $r_\odot = 8.33$~kpc and $r_s \simeq 14.0$~kpc.
We can adjust these limits using our parameters by rescaling them with the ratio of the $J$-factors for the two sets of parameters, 
both computed over the region of observation of the Inner Galaxy Survey,
$0.5^\circ < \theta < 3^\circ$ with $|b| > 0.3^\circ$, which in this case
gives $J^{\mathrm{H.E.S.S.}}/J^{\mathrm{our}} \simeq 0.60$.

Since the expected gamma-ray flux is proportional to 
$\langle\sigma v\rangle J$, a cored NFW profile, having a smaller central DM 
density and hence a smaller $J$-factor, leads to a weaker limit:


\begin{equation}
\label{eq:cored_rescaling}
    \langle\sigma v\rangle_{\mathrm{lim}}^{\mathrm{cored}}
    \simeq
    \langle\sigma v\rangle_{\mathrm{lim}}^{\mathrm{NFW}}
    \frac{J_{\mathrm{NFW}}}{J_{\mathrm{cored}}}
    >
    \langle\sigma v\rangle_{\mathrm{lim}}^{\mathrm{NFW}},
\end{equation}


where $J_{\mathrm{NFW}}/J_{\mathrm{cored}} \simeq 4.0$ for a flat core of 
radius $r_c = 1$~kpc.

The resulting limits are shown in \cref{Cross_section} for the NFW and the 
cored NFW profiles, together with the thermal relic cross section of our 
model. We focus on the NFW and cored NFW profiles since the H.E.S.S. experimental limits we employ
are provided for the NFW profile.


\begin{figure}[H]
\centering
\includegraphics[width=0.93\textwidth]{Sigmav_limits.png}
\caption{\emph{95\% C.L. upper limits on the annihilation cross section in the
$W^{+}W^{-}$ channel from the H.E.S.S. Inner Galaxy Survey, rescaled to our
halo parameters. We adopt the NFW and the cored NFW profiles. The dashed line is the cross section that reproduces the
observed relic abundance in our model.}}
\label{Cross_section}
\end{figure}



It is useful to discuss the relation between the number of photons detected and
the DM mass in order to understand the shape of the limit curves. 
Recalling \cref{eq:AnnihilationFlux_general}, we notice that the 
$\gamma$-ray flux scales as

\begin{equation}
\label{eq:flux_scaling}
\Phi \propto \frac{\langle\sigma v\rangle}{m_{\mathrm{DM}}^2}\, J
\int_{E_{\mathrm{thr}}}^{m_{\mathrm{DM}}} \frac{dN_{f\to\gamma}}{dE}\, dE,
\end{equation}

where $E_{\mathrm{thr}}$ is the energy threshold of the telescope and the 
spectrum extends up to $E \simeq m_{\mathrm{DM}}$.
The upper limit $\langle\sigma v\rangle_{\mathrm{lim}}$ is inversely proportional to the
number of detectable photons per annihilation, hence we can distinguish three different regimes.

At low masses, most of the photons produced fall below 
$E_{\mathrm{thr}}$ and are lost, therefore the limit is weak and the curve rises for small 
$m_{\mathrm{DM}}$.

At intermediate masses, of the order 
of a TeV, the photon spectrum falls in the energy range where H.E.S.S. is 
most sensitive, and the limit reaches its strongest value, which corresponds
to the minimum of the curve.

At high masses, the telescopes are still able to detect most of the photons, but the annihilation rate is suppressed by the factor 
$1/m_{\mathrm{DM}}^2$ of \cref{eq:flux_scaling}: at fixed mass density, the number density of DM 
particles decreases as $\rho/m_{\mathrm{DM}}$, and two particles are needed 
for each annihilation, therefore the upper limit curve rises.

We can compare the cross section predicted by our model with the experimental upper limit and exclude the DM mass values that give a cross section larger than the upper limit.
Since our model predicts a constant cross section after imposing the relic abundance, the condition of exclusion becomes

\begin{equation}
    2.2\times10^{-26}~\mathrm{cm^3/s}
    >
    \langle\sigma v\rangle_{\mathrm{lim}}(m_{\mathrm{DM}}) .
\end{equation}
 
From \cref{Cross_section} we can evaluate this condition: 
if the upper limit lies below the thermal relic value, the corresponding DM mass is excluded. 
If it lies above the thermal relic value, the model is still allowed.

\Cref{Cross_section} also shows how the upper limit changes when different DM density profiles are assumed. 

For the NFW profile, the thermal relic cross section lies below the upper 
limit at all masses, so no value of $m_{\mathrm{DM}}$ is excluded within 
this approximation. Since the cored NFW yields an even weaker limit, as 
shown in \cref{eq:cored_rescaling}, the same conclusion holds for the cored 
profile: the thermal relic model is not excluded in either case.


It is important to notice that in this case we are assuming that our relic model is correct and then excluding the different DM masses and density profiles.
Strictly speaking, the rigorous interpretation is that a point $(m_{\mathrm{DM}}, \langle\sigma v\rangle)$ above the limit curve is excluded in the sense that the three hypotheses, which are the DM profile, mass and cross section, are simultaneously incompatible.
In our work we fix the cross section since it is a robust theoretical prediction,
while the DM density profile is affected by numerous uncertainties, which will be discussed in detail in the next chapter.
